% !TEX program = platex
% 数学ノート集・共通テンプレート v2（積分論2026年度第1部の版面を参考）
% A5判。章内で定義・定理・例・注意等の番号を共通化。
% コンパイル： platex -interaction=nonstopmode math_note_template_v2.tex (2回)
%              dvipdfmx math_note_template_v2.dvi
\documentclass[a5paper,10pt,onecolumn]{jsarticle}
\usepackage[dvipdfmx,top=22mm,bottom=20mm,inner=19mm,outer=18mm,headheight=13pt,headsep=6mm]{geometry}
\usepackage{amsmath,amssymb,amsthm,mathtools}
\usepackage{enumitem,fancyhdr,titlesec}
\usepackage[dvipdfmx,bookmarks=true,colorlinks=false]{hyperref}
\usepackage{pxjahyper}

% ----- 差し替える文書情報 -----
\newcommand{\NoteTitle}{変分原理とLagrange力学}
\newcommand{\NoteSeries}{数学ノート集}

% ----- 文字・組版 -----
% 参考PDFの質素な明朝＋Computer Modernの印象を引き継ぎながら，
% A5判で窮屈にならないよう本文の行間を少し確保する．
\linespread{1.13}
\setlength{\parindent}{1zw}
\setlength{\parskip}{0pt}
\setlength{\abovedisplayskip}{6pt plus 2pt minus 2pt}
\setlength{\belowdisplayskip}{6pt plus 2pt minus 2pt}
\setlength{\abovedisplayshortskip}{3pt plus 2pt minus 1pt}
\setlength{\belowdisplayshortskip}{4pt plus 2pt minus 1pt}
\setlist[enumerate]{topsep=3pt,itemsep=1pt,parsep=0pt,leftmargin=2.4zw}
\titleformat{\section}{\normalfont\normalsize\bfseries}{\thesection}{0.8zw}{}
\titlespacing*{\section}{0pt}{2.0ex plus .4ex}{.9ex}
\titleformat{\subsection}{\normalfont\normalsize\bfseries}{\thesubsection}{0.7zw}{}
\titlespacing*{\subsection}{0pt}{1.5ex plus .3ex}{.6ex}
\raggedbottom

% ----- 柱とページ番号：ページ上部、罫線なし -----
\fancypagestyle{mathnote}{%
  \fancyhf{}%
  \fancyhead[L]{\footnotesize \NoteSeries\ （\NoteTitle）}%
  \fancyhead[R]{\footnotesize \thepage}%
  \renewcommand{\headrulewidth}{0pt}%
  \renewcommand{\footrulewidth}{0pt}%
}
\pagestyle{mathnote}

% ----- 節内通し番号の定義・定理群 -----
\newtheoremstyle{mathnotetheorem}
 {4.5pt}{4.5pt}{\normalfont}{}%
 {\bfseries}{}{0.65zw}%
 {\thmname{#1}\thmnumber{ #2}\thmnote{（#3）}}
\theoremstyle{mathnotetheorem}
\newtheorem{theorem}{定理}[section]
\newtheorem{definition}[theorem]{定義}
\newtheorem{lemma}[theorem]{補題}
\newtheorem{proposition}[theorem]{命題}
\newtheorem{corollary}[theorem]{系}
\newtheorem{example}[theorem]{例}
\newtheorem{remark}[theorem]{注意}
\renewcommand{\proofname}{証明}
\renewcommand{\qedsymbol}{$\square$}
\makeatletter
\renewenvironment{proof}[1][\proofname]{%
  \par\pushQED{\qed}\normalfont\topsep4\p@\@plus2\p@\relax
  \trivlist\item[\hskip\labelsep\bfseries #1]\ignorespaces
}{\popQED\endtrivlist\@endpefalse}
\makeatother

% ----- 使うノートで適宜拡張する記号 -----
\newcommand{\R}{\mathbb{R}}
\newcommand{\norm}[1]{\lVert #1\rVert}
\newcommand{\abs}[1]{\lvert #1\rvert}

\begin{document}
\section*{到達目標と前提}
作用汎関数からEuler--Lagrange方程式を導き，
具体的な境界条件を満たす停留曲線を構成できることを
主目標とする．保存則は得られた運動を検算する手段となる．
前提は微積分・常微分方程式・線形代数の初歩である．


\par\noindent\textbf{記号と約束.}\quad
関数空間の添字 $c$ はコンパクトな台をもつことを表す．
\par

\section{作用汎関数}
本稿の出発点は，配置変数 $q(t)\in\mathbb R^d$ と
Lagrange関数
$L:\mathbb R^d\times\mathbb R^d\times[a,b]\to\mathbb R$
である．物理的なLagrange関数の選択そのものはモデルの仮定であり，
その後の微分方程式は数学的な帰結として導く．
\begin{definition}
$L$ を $C^2$ 級とし，$q\in C^2([a,b];\mathbb R^d)$ に対して
作用汎関数を
\[
 \mathcal S[q]=\int_a^bL(q(t),\dot q(t),t)\,dt
\]
と定める．固定端点を持つ変分とは
$q_\varepsilon=q+\varepsilon\eta$ で
$\eta\in C^2([a,b];\mathbb R^d)$，
$\eta(a)=\eta(b)=0$ としたものをいう．
\end{definition}
\begin{definition}
$q$ が作用の停留曲線であるとは，
任意の固定端点変分に対し
$\left.\frac d{d\varepsilon}\right|_{\varepsilon=0}
 \mathcal S[q_\varepsilon]=0$ となることをいう．
\end{definition}
\begin{lemma}[滑らかな試験関数の存在]
任意の非空開区間 $J\subset(a,b)$ に対し，
$\eta\in C_c^\infty(J)$ で $\eta\ge0$ かつ
$\eta\not\equiv0$ となるものが存在する．
\end{lemma}
\begin{proof}
$\rho(t)=e^{-1/t}$（$t>0$），$\rho(t)=0$（$t\le0$）
と定める．$t>0$ での任意階導関数は
$1/t$ の多項式と $e^{-1/t}$ の積なので
$t\downarrow0$ で零に収束する．
従って $\rho$ は $\mathbb R$ 上の滑らかな非負関数である．
$[u,v]\subset J$ となる $u<v$ を取り，
$\eta(t)=\rho(t-u)\rho(v-t)$ と置けば，
$\eta$ は $[u,v]$ に台を持つ滑らかな非負関数であり，
$u<t<v$ では正である．
\end{proof}
\begin{lemma}[変分法の基本補題]\label{lem:fundamental}
$f\in C([a,b])$ が
$\int_a^bf(t)\eta(t)\,dt=0$ を
任意の $\eta\in C_c^\infty((a,b))$ について満たすなら，
$f(t)=0$ が全ての $a<t<b$ で成立する．
\end{lemma}
\begin{proof}
$f(t_0)>0$ となる $t_0\in(a,b)$ があると仮定する．
連続性から $t_0$ を含む小開区間 $J$ 上で
$f(t)\ge c>0$ とできる．
前補題の非負で零でない試験関数 $\eta$ を取ると
$\int f\eta\ge c\int\eta>0$ となり矛盾する．
$f(t_0)<0$ の場合にも $-f$ へ同じ議論を適用する．
従って全ての内部点で $f=0$ である．
\end{proof}

\section{Euler--Lagrange方程式}
\begin{theorem}[Euler--Lagrange方程式]\label{thm:el}
$L\in C^2$，$q\in C^2([a,b];\mathbb R^d)$ とする．
$q$ が固定端点変分に関して停留曲線であることと，
各 $i=1,\ldots,d$ に対して
\[
 {d\over dt}{\partial L\over\partial v_i}
 (q(t),\dot q(t),t)
 -{\partial L\over\partial q_i}
 (q(t),\dot q(t),t)=0
 \tag{EL}
\]
が $a<t<b$ で成立することは同値である．
\end{theorem}
\begin{proof}
$\eta\in C^2([a,b];\mathbb R^d)$，
$\eta(a)=\eta(b)=0$ を取る．
$L$ は $C^2$ なので，コンパクト区間上の連鎖律と
微分積分交換によって
\[
 \left.{d\over d\varepsilon}\right|_0\mathcal S[q+\varepsilon\eta]
 =\int_a^b
 \sum_{i=1}^d
 \left(
 {\partial L\over\partial q_i}\eta_i+
 {\partial L\over\partial v_i}\dot\eta_i
 \right)dt
\]
を得る．偏導関数は曲線 $(q,\dot q,t)$ 上で評価する．
$\eta_i(a)=\eta_i(b)=0$ なので，部分積分により
\[
 \left.{d\over d\varepsilon}\right|_0\mathcal S[q+\varepsilon\eta]
 =\int_a^b\sum_{i=1}^d
 \left({\partial L\over\partial q_i}
 -{d\over dt}{\partial L\over\partial v_i}\right)\eta_i\,dt.
 \tag{1}
\]
$q$ が停留曲線なら左辺は任意の $\eta$ に対して零である．
一つの成分 $\eta_i$ だけを任意の内部コンパクト台の試験関数とし，
残る成分を零と置けるから，
基本補題\ref{lem:fundamental}によって括弧内は零となる．
逆に (EL) が成立すれば (1) の右辺は零だから
$q$ は停留曲線である．
\end{proof}
\begin{remark}
停留曲線という条件は最小値を取ることを意味しない．
Euler--Lagrange方程式は固定端点問題の一階変分が
消えるための必要十分条件であり，
極小性には二階変分等の追加条件が必要である．
\end{remark}
\begin{example}[停留性と極小性の違い]\label{ex:nonminimum}
$T>\pi$ とし，$q(0)=q(T)=0$ を課した
\[
 \mathcal S[q]={1\over2}\int_0^T(\dot q^2-q^2)\,dt
\]
を考える．$q\equiv0$ は停留曲線であるが局所最小ではない．
\end{example}
\begin{proof}
$L=(v^2-q^2)/2$ のEuler--Lagrange方程式は
$\ddot q+q=0$ であり，$q\equiv0$ は解である．
$h(t)=\sin(\pi t/T)$ と置くと端点で零であり，
直接積分して
\[
 \mathcal S[\varepsilon h]
 ={\varepsilon^2T\over4}
 \left((\pi/T)^2-1\right)<0=\mathcal S[0]
\]
を得る．任意に小さい $\varepsilon\ne0$ で作用が小さくなるため
局所最小ではない．
\end{proof}
\section{運動方程式と具体例}
\begin{theorem}[Newton方程式の変分原理]
正定値対称行列 $M\in\mathbb R^{d\times d}$ と
$V\in C^2(\mathbb R^d)$ に対し
\[
 L(q,v)={1\over2}v^TMv-V(q)
\]
と置く．そのEuler--Lagrange方程式は
\[
 M\ddot q=-\nabla V(q)
\]
である．
\end{theorem}
\begin{proof}
$\partial L/\partial v=Mv$，
$\partial L/\partial q=-\nabla V(q)$ である．
$M$ は定数なので
$d(M\dot q)/dt=M\ddot q$．
これらを (EL) に代入すれば結論を得る．
\end{proof}
\begin{example}[調和振動子]
$L(q,v)=\tfrac12mv^2-\tfrac12kq^2$，
$m,k>0$ とすると，
全ての運動解は
\[
 q(t)=A\cos(\omega t)+B\sin(\omega t),\qquad
 \omega=\sqrt{k/m}
\]
の形である．
\end{example}
\begin{proof}
Euler--Lagrange方程式は
$m\ddot q+kq=0$ である．
表示した関数はこの方程式を満たす．
任意の初期位置と初期速度は
$A=q(0)$，$B=\dot q(0)/\omega$ を選べば実現される．
線形常微分方程式の初期値解は一意なので，
これらが全ての解である．
\end{proof}

\section{保存量とNoetherの定理}
\begin{definition}
一般化運動量とエネルギー関数を
\[
 p_i={\partial L\over\partial v_i},\qquad
 E(q,v,t)=\sum_{i=1}^dp_iv_i-L(q,v,t)
\]
によって定める．
\end{definition}
\begin{theorem}[時間並進不変性とエネルギー保存]\label{thm:energy}
$q$ がEuler--Lagrange方程式の解なら
\[
 {d\over dt}E(q(t),\dot q(t),t)
 =-{\partial L\over\partial t}(q(t),\dot q(t),t)
\]
が成立する．特に $L$ が $t$ に陽に依存しなければ
$E$ は運動に沿って一定である．
\end{theorem}
\begin{proof}
連鎖律を用いて
\[
 \begin{aligned}
 {dE\over dt}
 &=\sum_i(\dot p_i\dot q_i+p_i\ddot q_i)
 -\sum_i\left(
 {\partial L\over\partial q_i}\dot q_i+
 {\partial L\over\partial v_i}\ddot q_i
 \right)-{\partial L\over\partial t}\\
 &=\sum_i
 \left(\dot p_i-{\partial L\over\partial q_i}\right)\dot q_i
 -{\partial L\over\partial t}.
 \end{aligned}
\]
Euler--Lagrange方程式から括弧内は零なので，
表示した等式を得る．
\end{proof}
\begin{definition}
配置空間の局所一パラメータ変換群
$\Phi_s:\mathbb R^d\to\mathbb R^d$ を考える．
$\Phi_0=\mathrm{id}$，
$\Phi_{s+t}=\Phi_s\circ\Phi_t$ とし，
\[
 L(\Phi_s(q),D\Phi_s(q)v,t)=L(q,v,t)
\]
が成立するとき，$L$ は $\Phi_s$ に関して不変という．
その無限小生成ベクトル場を
$\xi(q)=\left.\partial_s\Phi_s(q)\right|_{s=0}$ とする．
\end{definition}
\begin{theorem}[Noetherの第一定理：点変換の場合]\label{thm:noether}
$L$ が上の一パラメータ変換群に関して不変とする．
Euler--Lagrange方程式の任意の解 $q(t)$ に沿って
\[
 J(t)=\sum_{i=1}^d
 {\partial L\over\partial v_i}(q,\dot q,t)\,
 \xi^i(q(t))
\]
は一定である．
\end{theorem}
\begin{proof}
不変性の式を $s=0$ で微分すると
\[
 \sum_i{\partial L\over\partial q_i}\xi^i(q)
 +\sum_{i,j}{\partial L\over\partial v_i}
 {\partial\xi^i\over\partial q_j}(q)v_j=0
 \tag{2}
\]
を得る．
一方解 $q(t)$ に沿って
\[
 {dJ\over dt}
 =\sum_i\dot p_i\xi^i(q)
 +\sum_{i,j}p_i{\partial\xi^i\over\partial q_j}(q)\dot q_j.
\]
Euler--Lagrange方程式は
$\dot p_i=\partial L/\partial q_i$ なので，
右辺は (2) の $v=\dot q$ に等しく，零となる．
従って $J$ は保存される．
\end{proof}
\begin{corollary}[空間並進対称性]
$L$ が固定ベクトル $a\in\mathbb R^d$ に沿う並進
$q\mapsto q+sa$ に関して不変なら，
運動量 $p\cdot a$ は保存される．
\end{corollary}
\begin{proof}
並進の無限小生成ベクトル場は定数場 $\xi(q)=a$ なので，
定理\ref{thm:noether}の保存量は $p\cdot a$ である．
\end{proof}
\begin{corollary}[回転対称性と角運動量]
$d=3$ とし，$L$ が固定した軸 $a$ の周りの
一パラメータ回転に関して不変なら，
$a\cdot(q\times p)$ は保存される．
\end{corollary}
\begin{proof}
回転の無限小生成ベクトル場は
$\xi(q)=a\times q$ である．
従ってNoetherの保存量は
$p\cdot(a\times q)=a\cdot(q\times p)$ となる．
\end{proof}

\section{Legendre変換とHamilton方程式}
\begin{definition}
$L$ が正則であるとは，速度変数に関するHessian
\[
 W(q,v,t)=\left(
 {\partial^2L\over\partial v_i\partial v_j}(q,v,t)
 \right)_{i,j}
\]
が可逆であることをいう．
\end{definition}
\begin{theorem}[局所Legendre変換]\label{thm:legendre}
$L\in C^2$ とし，ある $(q_0,v_0,t_0)$ で
$\det W(q_0,v_0,t_0)\ne0$ とする．
$p=\partial L/\partial v$ の関係を
$(q_0,v_0,t_0)$ の近傍で
$v=v(q,p,t)$ と一意に解ける．
このとき
\[
 H(q,p,t)=p\cdot v(q,p,t)-L(q,v(q,p,t),t)
\]
と置けば
\[
 {\partial H\over\partial p_i}=v_i,\qquad
 {\partial H\over\partial q_i}
 =-{\partial L\over\partial q_i},\qquad
 {\partial H\over\partial t}=-{\partial L\over\partial t}
 \tag{3}
\]
が成立する．
\end{theorem}
\begin{proof}
$F(q,v,t,p)=\partial L/\partial v(q,v,t)-p$ と置く．
$v$ に関する微分は $W$ なので，陰関数定理から
$v=v(q,p,t)$ を近傍で一意に解ける．
$H$ を $p_i$ で微分すると
\[
 {\partial H\over\partial p_i}
 =v_i+\sum_jp_j{\partial v_j\over\partial p_i}
 -\sum_j{\partial L\over\partial v_j}
  {\partial v_j\over\partial p_i}=v_i
\]
である．最後の等号は $p_j=\partial L/\partial v_j$ による．
$q_i,t$ についても連鎖律による速度微分項が相殺するため，
残りの二式を得る．
\end{proof}
\begin{corollary}[Hamiltonの正準方程式]
正則な $L$ のEuler--Lagrange解について
$p_i=\partial L/\partial v_i$ と置く．
その運動は局所的に
\[
 \dot q^i={\partial H\over\partial p_i},\qquad
 \dot p_i=-{\partial H\over\partial q^i}
\]
で記述される．
\end{corollary}
\begin{proof}
第1式は定理\ref{thm:legendre}の
$H_{p_i}=v_i$ と $v_i=\dot q^i$ による．
第2式はEuler--Lagrange方程式
$\dot p_i=L_{q^i}$ と $H_{q^i}=-L_{q^i}$ から従う．
\end{proof}
\section{Riemann幾何学への接続}
\begin{theorem}[測地線の変分原理]
Riemann多様体 $(M,g)$ の局所座標において
\[
 L(x,v)={1\over2}\sum_{i,j}g_{ij}(x)v^iv^j
\]
と置く．そのEuler--Lagrange方程式は
\[
 \ddot x^\ell+\sum_{i,j}\Gamma^\ell_{ij}(x)
 \dot x^i\dot x^j=0
\]
であり，Levi-Civita接続に関する測地線方程式に一致する．
\end{theorem}
\begin{proof}
$g_{ij}=g_{ji}$ なので
\[
 {\partial L\over\partial v^k}=\sum_jg_{kj}v^j,\qquad
 {\partial L\over\partial x^k}
 =\tfrac12\sum_{i,j}\partial_k g_{ij}v^iv^j.
\]
Euler--Lagrange方程式へ代入すると
\[
 \sum_jg_{kj}\ddot x^j+
 \sum_{i,j}\left(
 \tfrac12\partial_i g_{kj}
 +\tfrac12\partial_jg_{ki}
 -\tfrac12\partial_kg_{ij}\right)
 \dot x^i\dot x^j=0.
\]
両辺に逆行列 $g^{\ell k}$ を掛ければ
\[
 \ddot x^\ell+\tfrac12\sum_{i,j,k}g^{\ell k}
 (\partial_i g_{kj}+\partial_jg_{ki}-\partial_kg_{ij})
 \dot x^i\dot x^j=0
\]
を得る．括弧の係数はChristoffel記号
$\Gamma^\ell_{ij}$ なので結論が従う．
\end{proof}
\begin{remark}
この変分原理は座標に依存しない運動方程式を
接続の言葉で記述できることを示す．
大域的な極小測地線の存在には，
Euler--Lagrange方程式だけでなく
完備性などの条件が必要である．
\end{remark}

\section{端点が与えられた振動子を実際に解く}
\begin{proposition}[調和振動子の固定端点問題]
$m,k>0$，$\omega=\sqrt{k/m}$ として
\[
 L(q,\dot q)=\frac m2\dot q^2-\frac k2q^2
\]
を考える．$q(0)=a$，$q(T)=b$ を満たす
作用の停留曲線について，$\sin(\omega T)\ne0$ なら
ただ一つ存在し，
\[
 q(t)=
 \frac{a\sin(\omega(T-t))+b\sin(\omega t)}
      {\sin(\omega T)}
\]
で与えられる．一方 $T=\pi/\omega$ の場合，
$b=-a$ なら無限個存在し，$b\ne-a$ なら存在しない．
\end{proposition}
\begin{proof}
Euler--Lagrange方程式は
$m\ddot q+kq=0$ すなわち
$\ddot q+\omega^2q=0$ である．その一般解は
$q(t)=A\cos(\omega t)+B\sin(\omega t)$ である．
初期条件から $A=a$，
終端条件から
\[
 b=a\cos(\omega T)+B\sin(\omega T).
\]
$\sin(\omega T)\ne0$ なら $B$ が一意に定まり，
三角関数の加法公式で書き直すと表示式となる．
$T=\pi/\omega$ なら
$\cos(\omega T)=-1$，$\sin(\omega T)=0$
だから終端条件は $b=-a$ に限られ，
成立するとき $B$ は任意である．
これらの曲線が停留曲線となることは
Euler--Lagrange方程式の必要十分性から従う．
\end{proof}
\begin{example}[停留曲線とエネルギー]
$m=k=1$，$T=\pi/2$，$a=0$，$b=1$
とすると $q(t)=\sin t$ となる．
エネルギー
\[
 E(t)=\frac12\dot q(t)^2+\frac12q(t)^2
     =\frac12\cos^2t+\frac12\sin^2t
     =\frac12
\]
は一定である．一般の解についても
$E'(t)=\dot q(t)(m\ddot q(t)+kq(t))=0$
である．このように保存量は境界値解を検算できる．
\end{example}
\section*{到達確認問題と解答}
\noindent\textbf{問題．}
$m=k=1$，$T=\pi/2$，$a=1$，$b=0$
の停留曲線を求め，エネルギーを計算せよ．

\medskip\noindent\textbf{解答．}
前命題から $q(t)=\cos t$．
$\dot q=-\sin t$ なのでエネルギーは
$E(t)=\frac12(\sin^2t+\cos^2t)=1/2$ である．
\section*{限界}
停留曲線であることは作用を最小化することを
無条件には意味しない．特に $\sin(\omega T)=0$ の
境界値問題では解の存在と一意性さえ変化するため，
変分問題では時間区間と境界条件が本質的である．

\end{document}
